\section{Manuel de montage du démonstrateur}
\label{chap:manuel_de_montage}
Le montage du robot est réalisé par sous-ensembles afin de faciliter
l'assemblage. 

La procédure est divisée en cinq étapes
principales : le montage de la base, le montage de la partie épaule, le montage de la partie coude, le montage
de la partie poignet, puis l'assemblage final des différents sous-ensembles.
\subsection{Montage de la base}
La première étape consiste à préparer la structure porteuse du robot. Les profilés \textit{item} sont montés sur la plaque de base.

\begin{enumerate}
    \item Positionner les profilés \textit{item} sur la plaque de base selon le plan d'assemblage.
    
    \item Fixer les profilés à l'aide des vis adaptés.
    \item Monter les poignées de transport sur les profilés.
\end{enumerate} 
\subsection{Montage de la partie épaule}

La partie épaule constitue la liaison entre la base et la partie coude. Elle permet la rotation du premier bras autour de l'axe vertical et regroupe les fonctions de guidage, de transmission et de motorisation.

\begin{enumerate}

    \item Placer la poulie équipée de son palier lisse sur l'axe de rotation.

    \item Installer les circlips permettant de maintenir axialement l'ensemble poulie--palier lisse.

    \item Introduire la courroie autour de la poulie pendant cette étape, afin de permettre la transmission du mouvement avec le second moteur.

    \item Positionner le bras sur l'axe de rotation.

    \item Mettre en place les roulements à billes sur l'arbre.

    \item Installer les circlips assurant le maintien axial des roulements.

    \item Placer l'accouplement flexible sur l'extrémité supérieure de l'arbre.

    \item Positionner le moteur d'entraînement du bras en l'alignant avec l'accouplement flexible.

    \item Fixer le moteur d'entraînement du bras sur son support.

    \item Installer le second moteur, pour entraîner le système de poulies et de courroie.


\end{enumerate}

\subsection{Montage de la partie coude}
La partie coude constitue une articulation intermédiaire du robot. Elle doit permettre la liaison entre les bras.
\begin{enumerate}
    \item Monter le premier circlips sur l'axe afin de définir une butée axiale.
    \item Insérer le roulement sur l'axe.
    \item Monter le second circlips afin de maintenir le roulement en position.
    \item Glisser l'arbre dans l'alésage.
    \item Positionner la poulie dans le palier lisse.
    \item Glisser l'ensemble poulie--palier sur l'axe.
    \item Maintenir l'ensemble axialement entre deux circlips.
\end{enumerate}

\subsection{Montage de la partie poignet}

La partie poignet permet la liaison avec le support du préhenseur. 

\begin{enumerate}
    \item Monter le roulement autour du support du préhenseur.
    \item Positionner les circlips afin de maintenir le roulement axialement.
    \item Monter le bras supérieur sur l'ensemble support--roulement.
    \item Mettre en place la pièce de serrage du roulement.
    \item Visser la pièce de serrage afin de bloquer correctement le roulement.
    \item Monter le bras inférieur sur le méplat du support de préhenseur.
    \item Visser le bras inférieur.
    \item Monter l'assemblage \textit{ball spline}.
    \item Monter les poulies supérieure et inférieure.
\end{enumerate}

\subsection{Assemblage final des sous-ensembles}
Une fois les sous-ensembles préparés, les différentes parties du robot sont assemblées entre elles. La liaison principale est réalisée au niveau du coude.
\begin{enumerate}
    \item Positionner la partie coude entre les bras correspondants.
    \item Insérer la vis de liaison.
    \item Serrer la vis afin d'assurer le maintien mécanique des bras.
    \item Contrôler l'absence de frottements excessifs, de jeu anormal ou de blocage.
\end{enumerate}